\chapter{The Renderer Contract}
\label{ch:renderer-contract}
\label{ch:backend-architecture}

Every renderer family implements the same internal C\texttt{++} contract, but the contract is
deliberately permissive. Some methods are pure virtual; others return null, throw, forward to a
reduced operation, or do nothing. Consequently, reaching an interface method is not evidence
that the selected family performed the requested work. This chapter defines that shared boundary.
Chapters~\ref{ch:renderer-selection-identity}--\ref{ch:diagnostic-renderers} own family-specific
mechanisms and results; Appendix~\ref{ch:feature-matrix} owns the complete family inventory.

\begin{evidencenote}
Interface declarations, factory routes, and shared public forwarding are source-proven at
\CnaRevisionShort. Several cross-family observations below come from a retained fourteen-product
audit. Those rows are historically recorded where they were not rerun for this edition. The
removed \texttt{ASCII} identity appears only when its historical result explains the current
\cnaclass{AsciiPostProcessEffect}; omission of a newer family means ``not measured by that
cohort,'' not ``unsupported.''
\end{evidencenote}

\section{Contract shape and default behavior}

The contract is one header:
\path{modules/graphics/include/CNA/Internal/Renderers/Common/IGraphicsRenderer.hpp}.
Eleven interfaces divide resource operations from device-wide dispatch:

\begin{center}
\small
\begin{tabular}{@{}p{0.32\textwidth}p{0.60\textwidth}@{}}
\toprule
\textbf{Interface} & \textbf{Responsibility} \\
\midrule
\texttt{IVertexBufferRenderer} & vertex upload, update hints, and declaration \\
\texttt{IIndexBufferRenderer} & 16- and 32-bit index upload \\
\texttt{ITextureRenderer} & 2D texture data, interop, and optional CPU-shadow behavior \\
\texttt{ITexture3DRenderer} / \texttt{ITextureCubeRenderer} & volume and cube storage \\
\texttt{IRenderTargetRenderer} / \texttt{IRenderTargetCubeRenderer} & target binding, applied MSAA, and attachment reporting \\
\texttt{IEffectRenderer} & renderer-specific program compilation, binding, uniforms, and textures \\
\texttt{IOcclusionQueryRenderer} & query begin/end/completion/sample count \\
\texttt{ISpriteBatchRenderer} & batch state and the three renderer draw forms \\
\texttt{IGraphicsRenderer} & device lifecycle, factories, state, targets, draws, presentation, readback, and capability queries \\
\bottomrule
\end{tabular}
\end{center}

The default bodies form five distinct contracts:

\begin{center}
\small
\begin{tabular}{@{}p{0.18\textwidth}p{0.31\textwidth}p{0.42\textwidth}@{}}
\toprule
\textbf{Default} & \textbf{Representative routes} & \textbf{Caller-visible consequence} \\
\midrule
Pure virtual & required 2D texture, SpriteBatch, vertex-buffer, and index-buffer factories & A family must implement the route to compile. \\
Null factory & volume/cube textures, render targets, effects, and occlusion queries & Public construction may still produce an object with no private implementation. \\
No-op & state setters, SpriteBatch options, several uniform setters & Stored public state or a successful call may have no native effect. \\
Reduced fallback & effect-aware draws may call a colored draw & Work can complete after effect parameters were discarded. \\
Throw/refusal & backbuffer readback, instancing, selected upload paths & Failure is explicit, although exception type and timing vary. \\
\bottomrule
\end{tabular}
\end{center}

\begin{contractnote}
Treat every optional method as an independent capability. A non-null renderer object, a true
capability answer, and one successful draw establish different facts. Appendix~\ref{ch:verification-tiers}
defines the evidence labels used to keep those facts separate.
\end{contractnote}

\subsection{Resource ownership and device tracking}
\label{sec:backend-resource-tracking}

A public \cnaclass{GraphicsResource} owns its renderer handle, normally through
\texttt{unique\_ptr}; \cnaclass{Texture2D} uses \texttt{shared\_ptr}. The device owns neither.
It records raw \cnaclass{GraphicsResource} addresses so it can dispose registered resources
before destroying the renderer. \texttt{GraphicsDevice::Dispose()} moves and clears that list,
calls each resource's virtual \texttt{Dispose()}, raises \texttt{Disposing}, and then releases
the renderer. Re-entrant deregistration therefore does not invalidate the iteration.

The registry follows object addresses, not renderer-handle ownership. The
\cnaclass{GraphicsResource} copy and default move constructors carry the device pointer without
registering the destination or replacing the source address. A \cnaclass{Texture2D} copy has its
own disposed flag and shares the private handle; moving buffers, textures, or render targets can
leave the device tracking the moved-from value. If a surviving untracked object outlives device
disposal, its renderer resource may be released after the native context.

\begin{limitationnote}
Keep live graphics resources address-stable and device-scoped. Prefer non-relocating ownership,
do not move them across devices, and dispose every \cnaclass{Texture2D} alias before the device.
The pin has no located test for copy/move followed by device-first teardown. The sibling
\path{docs/graphics-resource-lifetime.md} is stale where it says destructor cleanup suppresses
\texttt{ResourceDestroyed}; the current implementation still calls that device hook.
\end{limitationnote}

Chapter~\ref{ch:graphicsdevice}, \S\ref{sec:resource-lifetime}, covers event ordering and the
remaining constructor/destructor gaps. Figure~\ref{fig:resource-lifetime} summarizes the
ownership relation.

\subsection{Null factories and delayed failure}
\label{sec:backend-null-factory-outcomes}

Null factory results do not share one public failure policy:

\begin{center}
\small
\begin{tabular}{@{}p{0.22\textwidth}p{0.31\textwidth}p{0.38\textwidth}@{}}
\toprule
\textbf{Private result} & \textbf{Public result} & \textbf{First use at the pin} \\
\midrule
Null occlusion query & Constructed \cnaclass{OcclusionQuery} & Begin/end do nothing; completion is false and count is zero. \\
Null volume/cube texture & Policy depends on the public wrapper & An honest \texttt{Texture3D} capability answer rejects construction; an optimistic answer can still produce a null object. Cube and volume transfers now throw instead of fabricating or discarding content. \\
Null 2D/cube target & Constructed target with no cached target renderer & Sampling and transfer operations reject missing storage; target binding throws before public binding state changes. \\
Null effect & Renderer-dependent public failure & Chapter~\ref{ch:shader-fx-gap} distinguishes null, invalid, throwing, and working routes. \\
\bottomrule
\end{tabular}
\end{center}

Software supplies 2D targets and cube-texture storage but still inherits the null factories for
volume textures, cube targets, and occlusion queries; several 2D families inherit a wider subset.
D3D9 can also return null according to device capabilities. A cube or volume texture passed to
\cnaclass{ShaderEffect} remains a sharp edge because its CNAEXT renderer accessor assumes that the
private owner exists. Construction alone therefore does not establish usable storage. No
individual capability value covers every texture and target factory.

The framework needs one consistent source-level policy---early
\cnaclass{NotSupportedException}, or explicit nullable availability plus guards throughout.
This book records the current behavior and does not choose or implement that policy.

\subsection{Borrowed native handles}
\label{sec:backend-native-handle-boundary}

\texttt{GetWindowInternal()} and \texttt{GetRendererInternal()} are pure virtual because the
shared game path needs a window, even though not every family has one. The returned pointers are
borrowed:

\begin{itemize}
  \item ordinary windowed families return their current \texttt{SDL\_Window*};
  \item \texttt{HEADLESS}, \texttt{SOFTWARE}, and D3D12 \texttt{HeadlessEXT} return null;
  \item only \texttt{SDL\_RENDERER} returns its own \texttt{SDL\_Renderer*}; other families
        normally return null for that method;
  \item \cnaclass{GameWindow::GetNativeSdlWindowEXT()} and the integer-shaped XNA
        \texttt{Handle} property expose the same window without transferring ownership.
\end{itemize}

Do not call \texttt{SDL\_DestroyWindow} on either result or cache it across device disposal.
Explicitly disposing the game-owned device does not invalidate the separate
\cnaclass{GameWindow} wrapper, so later use of that wrapper can reach a destroyed SDL object.

\texttt{PresentationParameters::DeviceWindowHandle} is the inverse borrowed route. A nonzero
value is reinterpreted as \texttt{SDL\_Window*}; CNA does not validate or destroy it. The pinned
implementation publishes an attached window to the global mouse and text-input services and, at
teardown, clears either handle only if it still names that window. An embedding host must keep the
window alive through renderer teardown. The global slots still have last-writer semantics, so two
live devices can replace one another's input target; the located tests do not cover that
multi-device lifecycle.

\section{Presentation and input are family contracts}

The public manager stores one presentation request. The selected family determines whether that
request changes a compositor, a swapchain, a CPU buffer, only reported state, or nothing.

\subsection{Input-coordinate routing}
\label{sec:backend-input-coordinate-matrix}

Window-to-logical input uses three tiers: an associated SDL renderer; otherwise a static
\texttt{SDL\_Window*} to \texttt{IGraphicsRenderer*} registry; otherwise raw coordinates. The
inverse mouse-warp route uses the same ordering. The compact matrix below names the current
behavioral classes; family chapters retain the implementation detail.

\begin{center}
\small
\begin{tabular}{@{}p{0.23\textwidth}p{0.31\textwidth}p{0.37\textwidth}@{}}
\toprule
\textbf{Family group} & \textbf{Route} & \textbf{Boundary} \\
\midrule
\texttt{SDL\_RENDERER} & SDL logical-presentation conversion & Implements all five modes, including offsets; public read-side and direct write-side tests exist. \\
EasyGL / Canvas & Static renderer registry & Fixed-height scaling is implemented; other stored modes lack matching offset/output behavior. \\
\texttt{SDL\_GPU} / \texttt{WEBGPU} & Static renderer registry & Five-mode math uses physical pixels while events and warps use window coordinates; high pixel density can mis-scale input. \\
\texttt{FREEDIRECT} & Associated private SDL renderer wins & Mapping appears after first present and is hard-coded to letterbox; direct CNA transform overrides are shadowed. \\
Vulkan, BGFX, Direct3D families & Raw fallback & Either presentation remains physical or the stored virtual request is not implemented as an output transform. \\
\texttt{SOFTWARE} / \texttt{HEADLESS} & No window route & Logical coordinates belong only to internal buffers or trace state. \\
\bottomrule
\end{tabular}
\end{center}

An immediate \texttt{Mouse::SetPosition}/\texttt{GetState} test is weak because the setter first
writes the requested logical value into \texttt{InputManager}; that round trip can pass even if
the physical warp is wrong. The strongest located public evidence injects physical SDL motion
into a scaled presentation and observes the converted logical point.

\subsection{Window-registry lifetime}
\label{sec:backend-window-registry-lifetime}

The second tier is one process-wide, unsynchronized
\texttt{unordered\_map<SDL\_Window*, IGraphicsRenderer*>}. Registration overwrites by key and
unregistration erases by window alone; neither operation checks ownership. EasyGL registers
before several fallible constructor steps, so an exception can leave a dangling value. Two
devices attached to the same borrowed window create another failure: the second replaces the
first, then destruction of the first removes the second's entry. Address reuse and concurrent
lookup add equivalent hazards.

This requires an ownership and thread-affinity decision in CNA. Suitable designs include an RAII
registration after fallible setup, owner-checked removal, and either rejecting duplicate devices
or defining multi-device routing. No located test covers constructor failure after registration,
duplicate attachment, window-address reuse, or concurrent access.

\subsection{Presentation modes}
\label{sec:backend-presentation-mode-matrix}

The CNAEXT \cnaclass{PresentationMode} property belongs to
\cnaclass{GraphicsDeviceManager}. Direct device construction leaves the renderer-create argument
at its fixed-height default. The manager route applies the mode before reset so implementations
that derive a virtual width see the intended value.

\begin{center}
\small
\begin{tabular}{@{}p{0.23\textwidth}p{0.31\textwidth}p{0.37\textwidth}@{}}
\toprule
\textbf{Implementation group} & \textbf{Current behavior} & \textbf{Evidence scope} \\
\midrule
\texttt{SDL\_RENDERER} & Maps all five requests to SDL logical presentation; fixed height derives width. & Source and focused routing tests; no five-mode pixel corpus. \\
\texttt{SDL\_GPU} / \texttt{WEBGPU} & Computes native, fixed-height, stretch, letterbox, and overscan rectangles. & Complete local geometry; visual coverage is partial. \\
EasyGL / Canvas & Derives fixed-height width; other values retain dimensions without implementing bars or crop. & Source plus fixed-height coverage; browser output remains unverified. \\
Vulkan / BGFX / Direct3D & Stores, ignores, or partially reports the request while rendering to physical output. & Source and renderer-specific smoke evidence; not a portable five-mode contract. \\
Software / Headless & Resizes or simulates an internal grid. & No display compositor exists, so bar/crop semantics do not apply. \\
\texttt{FREEDIRECT} & Stores the enum; the presenter uses letterbox and has a later-resolution defect. & Exact pixels for selected paths, not five policies. \\
\bottomrule
\end{tabular}
\end{center}

Portable code should configure the manager before normal initialization and use
\texttt{NativeBackBuffer} when physical readback coordinates matter. An enum round trip is not
evidence that letterbox, crop, or stretch reached output.

\subsection{Swap interval}
\label{sec:backend-swap-interval-matrix}

Public conversion preserves \texttt{Immediate}=0, \texttt{One}=1, and \texttt{Two}=2, but
families interpret them differently. The manager's Boolean retrace property reaches only 0 or 1;
2 requires explicit \cnaclass{PresentationParameters}. Because \cnaclass{Game} constructs its
first renderer before derived-game configuration, a family that consumes only the constructor
argument can remain at the default interval.

\begin{center}
\small
\begin{tabular}{@{}p{0.24\textwidth}p{0.28\textwidth}p{0.39\textwidth}@{}}
\toprule
\textbf{Group} & \textbf{0 / 1 / 2 mapping} & \textbf{Qualification} \\
\midrule
EasyGL & Passes the exact integer at construction and reset. & Off/on is runtime-observed through GL state; interval 2 is not timing-verified. \\
Vulkan & Immediate-or-mailbox / FIFO / FIFO-relaxed-or-FIFO at construction. & Runtime hook is a no-op; FIFO-relaxed is adaptive tearing, not half-refresh pacing. \\
BGFX, WebGPU, SDL GPU, D3D11/12 & Collapse positive values to enabled VSync or FIFO. & 2 is equivalent to 1 despite some downstream APIs supporting wider intervals. \\
D3D9 & Maps all three to D3D9 intervals and resets. & Windowed D3D9 may reject 2; CNA does not preflight that native restriction. \\
SDL Renderer & Initial positive values collapse to 1; runtime attempts the integer, then falls back to 1. & Tests prove safe reset and pixels, not driver timing. \\
Canvas, free-direct, Software, Headless & No display-interval implementation at this layer. & Stored public state may still change. \\
\bottomrule
\end{tabular}
\end{center}

No located timing test distinguishes half refresh from ordinary VSync on any family.

\subsection{Backbuffer format, depth format, and fullscreen}
\label{sec:backend-presentation-format-matrix}

Reset first asks the renderer to normalize the requested color and depth formats, then stages the
result as public state. Window resizing and virtual-resolution application form a rollback block:
if either throws, CNA restores the prior presentation, adapter, touch dimensions, and window size
before rethrowing. MSAA, interval, and the renderer's presentation-format update occur afterward;
they are not covered by that rollback. Only D3D9 overrides the final format-update hook. A family
whose applied-format query inherits the identity default can therefore echo a request even when
its native attachments use a fixed format.

\begin{center}
\small
\begin{tabular}{@{}p{0.21\textwidth}p{0.34\textwidth}p{0.36\textwidth}@{}}
\toprule
\textbf{Group} & \textbf{Native attachments} & \textbf{Reset/fullscreen boundary} \\
\midrule
D3D9 & Maps requested color and exact None/D16/D24/D24S8 depth forms. & Performs a device reset; unsupported combinations can fail after public state changed. \\
D3D11/12 & Fixed RGBA8 and D24S8 when a windowed backbuffer exists. & SDL changes the window; DXGI exclusive fullscreen is not requested. \\
Vulkan / WebGPU / SDL GPU / BGFX & Selects surface/device formats independently of the XNA enums; depth is fixed or device-selected. & Resize/surface recreation follows the window, not the stored enum fidelity. \\
EasyGL & SDL/GL chooses the default framebuffer; its MSAA FBO is fixed RGBA8/depth24 without stencil. & Tests establish stored values and selected depth behavior, not format reconfiguration. \\
SDL Renderer / Canvas / free-direct & 2D color target with no default depth/stencil attachment. & Fullscreen is an SDL/window or browser concern. \\
Software / Headless & Fixed CPU color/depth arrays, or trace state without attachments. & Fullscreen is stored state only. \\
\bottomrule
\end{tabular}
\end{center}

An SDL fullscreen failure is cleared and ignored, so \texttt{IsFullScreen} is a request, not an
actual-state query. D3D9 has the strongest located native format evidence; most other tests prove
stored-field round trips or continued rendering.

\subsection{Backbuffer readback}
\label{sec:backend-readback-matrix}

The interface promises top-left RGBA8 data and throws by default. The public wrapper checks a
null destination, positive in-bounds rectangles against the stored backbuffer dimensions, and
\texttt{elementCount < width*height}. It does not reject a negative start index, account for
remaining capacity after that offset, or guard the signed area multiplication against overflow.

Renderer behavior falls into four observable classes:

\begin{center}
\small
\begin{tabular}{@{}p{0.25\textwidth}p{0.28\textwidth}p{0.38\textwidth}@{}}
\toprule
\textbf{Class} & \textbf{Families} & \textbf{Image returned} \\
\midrule
Active target & SDL Renderer, EasyGL, Software, free-direct, Canvas; historical ASCII & Bound target when present, otherwise the family backbuffer/logical surface. Canvas browser execution was not observed for this edition. \\
Physical/default backbuffer & Vulkan, BGFX, D3D9, D3D11, WebGPU & Ignores an active user target. Logical and physical dimensions can diverge on scaled/high-density output. \\
Synthetic & Headless & Fills with the last global clear color; draw output and source coordinates are not represented. \\
Explicit refusal & SDL GPU, D3D12 & Inherited throw, even though both have separate render-target readback mechanisms. \\
\bottomrule
\end{tabular}
\end{center}

There is no capability query for this distinction. A portable test must name the selected
family, source image, coordinate space, and oracle.

\section{Resource and draw defaults}

\subsection{Target, effect, and query interfaces}

\texttt{IRenderTargetRenderer::GetMultiSampleCount()} defaults to zero, while
\texttt{HasRealDepthBuffer}\allowbreak\texttt{(requested)} defaults to echoing the request. The first is conservative;
the second can claim an attachment that the family never created. SDL Renderer overrides the
depth answer to false. Applied sample count and actual attachment presence are therefore stronger
facts than constructor arguments.

\texttt{IOcclusionQueryRenderer} itself is pure virtual, but its factory is optional.
\texttt{IEffectRenderer} requires program compile/bind/status methods while many scalar, array,
and texture setters have no-op defaults. A valid program can consequently discard selected
parameters. Chapter~\ref{ch:shader-fx-gap} owns the renderer-specific matrix.

\subsection{Buffer updates and instancing}

Dynamic buffer constructors currently use the same factory as static buffers; the constructor's
\texttt{dynamic} flag is ignored. Update options still have meaning on EasyGL, SDL GPU, D3D9,
and D3D11, which map discard/no-overwrite to distinct native operations. Vulkan and BGFX inherit
the option-discarding default; WebGPU, D3D12, Software, and Headless override it without
preserving option identity. The options-taking public wrappers also fail to refresh the shared
CPU shadow, so a later public \texttt{GetData()} can throw or return stale bytes.

The declaration route is stricter: \texttt{SetVertexDeclaration()} is pure virtual, and raw
uploads send the complete declaration. A family may translate it, remember it for draw-time
validation, or reject it. Chapter~\ref{ch:vertex-streams-capabilities} defines the packed-stream
contract.

\label{sec:drawex-instancing-matrix}
Effect-aware \texttt{DrawPrimitivesEx} defaults can discard \cnaclass{GpuDrawParams} and enter a
colored draw. The exhaustive family audit found DIRECTX10 as the consequential live inherited
fallback: an effect request can become a plausible untextured colored result. Several families
override the method but still use a colored fallback for unmatched layouts. Appendix
\ref{ch:feature-matrix} classifies inherited, hybrid, and renderer-owned routes.

Instancing has a throwing default. Shared code now transports up to sixteen bindings, preserving
slot, renderer buffer, stride, vertex offset, instance frequency, and count. It rejects partially
overlapping semantic declarations and validates per-vertex and per-instance ranges before
submission. The broad portable subset is one vertex stream plus one instance stream; multiple
streams of either rate additionally require \texttt{MultiStreamVertexInput}. EasyGL profiles
with native support, Magnum, Vulkan, BGFX, WebGPU, and D3D11/12 have focused pixel evidence for
specific layouts. D3D9 has a native step-frequency route but weaker shared-oracle coverage.
Headless records counts, and SDL GPU inherits a default-true capability despite the throwing
callback; neither establishes working instancing.

\subsection{Texture updates and readback}
\label{sec:texture2d-update-matrix}

Every current 2D texture implementation overrides level-zero updates. The higher-mip default is
still reachable on SDL GPU and silently discards the upload; Software explicitly discards higher
levels, while SDL Renderer, Canvas, and free-direct throw. EasyGL, Vulkan, BGFX, WebGPU, and
D3D9/11/12 implement higher-level storage. GPU-sampled or native-subresource tests, not public
CPU-shadow round trips, establish those paths. WebGPU additionally regenerates the mip chain
after level-zero writes, which can overwrite previously authored higher levels.

\label{sec:texture-volume-readback-matrix}
Volume and cube readback now returns a completion Boolean. Public wrappers convert results only
after a true response; null/false raises \cnaclass{NotSupportedException} without modifying the
destination. EasyGL, Vulkan, BGFX (when transfer capabilities exist), WebGPU, SDL GPU,
D3D9/11/12, and LLGL have native transfer routes. Skia and selected Software cube paths use exact
CPU storage. Headless refuses instead of manufacturing transparent pixels. The shared fixtures
distinguish slices, faces, subrectangles, mips, and unchanged sentinel buffers.

\label{sec:rendertargetcube-default-contracts}
\cnaclass{RenderTargetCube} inherits public texture upload methods, but most renderer-target
objects explicitly refuse them; EasyGL and Skia are the notable storage implementations.
Target transitions are a separate concern. Families with deferred work must finalize the
outgoing face before another cube face, 2D target, MRT set, or backbuffer becomes active.
D3D9/11/12, BGFX, WebGPU, SDL GPU, LLGL, and Vulkan have renderer-owned transition mechanisms.
EasyGL retains per-face MSAA storage, but switching faces on the same cube can skip outgoing
finalization; callers that need freshly resolved/sampleable faces should unbind between them.

\section{Targets, clears, viewport, and scissor}

\subsection{Render-target identity and usage}
\label{sec:backend-rendertarget-usage-matrix}

Singular 2D and cube calls normalize into one descriptor list. Each descriptor names a 2D target
or cube face, dimensions, and applied sample count. The first descriptor controls viewport,
scissor, discard clearing, and usage policy. The shared layer rejects null/disposed targets,
nonzero 2D slices, dimension or sample-count mismatches, duplicate subresources, and unsupported
mixed/cube MRT shapes before native binding.

Three rules are cross-family:

\begin{enumerate}
  \item CNA has no FNA-style redundant-binding early return. Rebinding the same discard target
        clears it again and can repeat resolve/finalization work.
  \item \texttt{PreserveContents} and \texttt{PlatformContents} both pass the preserve Boolean;
        exact \texttt{DiscardContents} triggers the shared implicit clear.
  \item The discard mask includes color and each depth/stencil aspect the first descriptor
        reports as present. Backbuffer usage is outside this policy.
\end{enumerate}

Focused fixtures distinguish discard from preserve/platform for major GPU families, including
cube and depth/stencil cases. Software has CPU color, depth, and stencil storage; Headless reports
trace changes, not preserved pixels. WebGPU and Software explicitly reject more than one target.
D3D12 binds multiple RTVs, but current stock draws write target zero. A true MRT capability answer
must therefore be paired with a successful multi-target bind and an oracle for the intended
outputs.

\subsection{Clear selection and ordering}
\label{sec:backend-clear-matrix}

\cnaclass{GraphicsDevice} converts color once, masks unavailable target aspects, and dispatches
one of seven non-empty target/depth/stencil combinations. The significant remaining differences
are concise:

\begin{itemize}
  \item Vulkan, BGFX, WebGPU, and SDL GPU preserve ordered clears across deferred pass segments;
        focused pixel suites cover masks, chronology, and several target kinds.
  \item D3D11/12 issue native immediate clears. D3D9 further gates depth/stencil bits using its
        backbuffer format, even when a differently formatted target is active.
  \item EasyGL's color-only hook also includes the GL depth bit without forcing depth writes, so
        prior depth-mask state can change a nominal target-only clear.
  \item SDL Renderer, Canvas, and free-direct mask depth/stencil away on their public 2D paths;
        their direct 3D hooks may throw but are not reached by the masked public call.
  \item Software handles color, depth, and stencil in CPU arrays; Headless records trace identity
        without attachment semantics.
\end{itemize}

There is no capability value for aspect selectivity or same-pass ordering. Tests must observe the
requested aspect after intervening work; an unchanged color pixel alone does not prove that a
depth-only clear occurred.

\subsection{Viewport and scissor}
\label{sec:backend-viewport-scissor-matrix}

The device calls each renderer hook before committing the corresponding public property, and
scissor enable arrives separately through \cnaclass{RasterizerState}. A successful getter round
trip therefore establishes accepted bookkeeping, not rasterization.

Vulkan, BGFX, WebGPU, and SDL GPU capture viewport/scissor per deferred command and have focused
pixel evidence. EasyGL and D3D9/11 apply both through native state; EasyGL SpriteBatch resets the
viewport to the full destination and does not restore it. D3D12 applies custom viewport and depth
range but installs a full-target scissor on every draw. SDL Renderer applies a nonempty clip
rectangle regardless of \texttt{ScissorTestEnable} and has no viewport hook. Software applies the
rectangle and enable state to its 2D and 3D clip paths; its custom viewport is covered for
SpriteBatch, while the 3D NDC mapping still spans the full framebuffer. Canvas, free-direct, and
historical ASCII do not implement both semantics. Headless validates/traces but does not rasterize.

Target switches reset viewport and scissor to destination bounds. Construction establishes both
full rectangles. Present and a same-size reset preserve custom values; a logical-size or physical
default-viewport change resets both. No capability value reports either feature.

\section{Capability answers and operational evidence}
\label{sec:graphicscapability}

\cnaclass{GraphicsCapability} contains \texttt{ThreeD}, depth/stencil, MSAA, MRT, anisotropy,
wireframe, occlusion query, custom effects, 3D textures, multi-stream input, instancing, stencil,
and additive blending. The shared default returns true except for multi-stream input; stencil
uses a separate hook. This optimistic polarity makes an answer useful only when the family has
audited the corresponding operation.

Figure~\ref{fig:capability-evidence} shows the required progression from reported Boolean to
an observed result.

Confirmed contradictions and scope mismatches at the pin include:

\begin{center}
\small
\begin{tabular}{@{}p{0.23\textwidth}p{0.26\textwidth}p{0.42\textwidth}@{}}
\toprule
\textbf{Capability} & \textbf{Affected families} & \textbf{Reachable behavior} \\
\midrule
MRT & Headless, Software, WebGPU & Trace-only behavior or explicit rejection, not a usable multi-attachment draw contract. \\
Occlusion query & WebGPU, SDL GPU, Software; Headless is synthetic & Null factory or non-GPU bookkeeping despite true reporting. \\
Custom effects & BGFX, WebGPU & Always-invalid object, null factory, or a later throwing route. \\
Instancing & Headless, SDL GPU & Trace-only result or inherited throwing callback. \\
MSAA & Software; several device-dependent families are overbroad & Software implements 4x color MSAA for render targets but not its backbuffer; Vulkan/WebGPU/SDL GPU/D3D9 can negotiate a request down while still reporting true. \\
Anisotropy & SDL GPU, Software; D3D9 is overbroad & Stored/ignored value, discarded sampler state, or device limit not consulted by the query. \\
\bottomrule
\end{tabular}
\end{center}

For MSAA, the applied value written back to presentation parameters or the target's own
\texttt{MultiSampleCount} is stronger than the Boolean. For all capabilities, the defensible
sequence is query, factory, operation, engagement, then a suitable behavioral or pixel oracle.

\subsection{Recovery, device-loss, and markers}
\label{sec:backend-recovery-debug-matrix}

\texttt{SetContextRecoveryEnabled()} changes shared texture-shadow policy before invoking the
renderer hook. Disabling it permits a successful full upload to discard CPU pixels, so later
readback or a partial update can fail even when the family hook is empty. Re-enabling recovery
does not reconstruct discarded data. Call it after device construction but before loading
resources; the header's ``before device initialization'' wording cannot be literal because the
method belongs to an existing device.

F9 and F10 key-down events pass through input first, then call renderer loss/restore hooks.
EasyGL uses them for desktop or browser context-loss routes. D3D9 exposes a staged lost/reset
lifecycle with event and pixel coverage. Other audited families generally inherit no-ops.
\texttt{SetStringMarkerEXT()} is also a no-op by default; Vulkan can queue a debug-utils marker
when the extension function exists, while D3D9 throws. Successful submission is not proof that a
native marker was emitted.

\section{Selection and factory boundary}

Chapter~\ref{ch:renderer-selection-identity} is the canonical identity/family registry. At this
edition, one configure-time public identity maps to one of \RendererFamilyCount{} implementation
families; the five EasyGL identities share one family. Platform and dependency gates run before
the factory is compiled.

The selected module provides exactly one
\texttt{CreateGraphicsRenderer(const GraphicsRendererCreateArgs\&)} returning
\texttt{unique\_ptr<IGraphicsRenderer>}. Its create arguments include window, virtual
resolution, presentation mode, recovery policy, requested samples and interval, plus CNAEXT
backbuffer/depth formats, fullscreen, and profile. Field presence is not a requirement that every
family consume the field; the presentation tables above show the current differences.

\begin{historynote}
The original audit narrated individual task numbers, intermediate tests, and repaired defaults
throughout this chapter. The current contract now incorporates those fixes. History remains here
only where it explains a live hazard: permissive defaults allowed plausible reduced draws;
CPU-shadow round trips masked missing GPU uploads; and direct helper tests bypassed the public
coordinate route. The detailed dated campaign remains in the repository audit artifacts.
\end{historynote}

\section{Implications for users and implementers}

For application code:

\begin{itemize}
  \item choose a renderer by the operations and evidence your program needs, not by selector name;
  \item treat capability answers as preflight hints and verify applied values or results;
  \item keep resources address-stable and within device lifetime;
  \item distinguish logical presentation coordinates from physical readback pixels;
  \item expect custom effects, MRT, queries, instancing, and readback to require family-specific
        qualification.
\end{itemize}

For a renderer implementation, overriding a virtual method is only the beginning. A complete
claim for one operation needs a reachable factory, shared public routing, applied state, correct
target and lifetime behavior, renderer-engagement evidence, and a discriminating oracle. That is
the standard used by the family chapters that follow.
